\section{Problem Formulation and Benchmark Overview}
\label{sec:benchmark-overview}

\subsection{Step-level routing}

A multi-turn LLM agent produces a trajectory $\tau = (s_1, s_2, \ldots, s_N)$ of LLM calls. At step $i$, the agent holds a message history $M_i = [m_1, \ldots, m_k]$ (a \emph{prefix}: system prompt, user instruction, previous assistant messages, tool outputs, retrieval chunks), emits a new assistant message $a_i = \mathcal{R}(M_i)$, executes any tool calls in $a_i$ against an external environment to produce observation $o_i$, and appends $a_i, o_i$ to the history: $M_{i+1} = M_i \parallel [a_i, o_i]$. The agent halts when it submits or when a budget is exhausted.

\subsection{Conditional per-call routing}
\label{sec:conditional-per-call-routing}

A standard LLM router maps a query to a candidate model to optimize quality and
cost; recent surveys formalize this as selecting from a model set subject to a
cost budget \citep{varangotreille2025doing}. TwinRouterBench studies the
conditional per-call version: the query is the full router-visible prefix $x_i$
before the next LLM call, including system messages, user requests, prior
assistant messages, tool outputs, retrieval snippets, logs, and partial code
edits. A step-level router is a conditional decision rule
\[
\pi : x_i \mapsto t_i \in \mathcal T,
\]
where
\[
\mathcal T =
\{\texttt{low},\texttt{mid},\texttt{mid\_high},\texttt{high}\}
\]
is an ordered set of cost and capability tiers. Each tier maps to a pool of models
with similar capability and cost; the concrete model is decided downstream by a
tier-to-model mapping, keeping public rows free of vendor model IDs.

The ideal conditional target tier is the cheapest tier that can
correctly handle the current step given the prefix $x_i$.
Let $\mathcal M_t\subseteq\mathcal M$ denote the model pool for tier
$t$, and let $V_i(m;x_i)=1$ indicate that model $m$ produces a sufficient
response for the current step under prefix $x_i$. For one-shot workloads,
$V_i$ is the task-level success predicate. For multi-turn agentic workloads,
$V_i(m;x_i)=1$ means that model $m$ correctly solves the problem posed at
step $i$; when each step is solved correctly, the trajectory as a whole is
more likely to succeed. The ideal target is
\[
t_i^\star
=
\min\left\{
t\in\mathcal T:
\exists m\in\mathcal M_t
\ \text{such that}\
V_i(m;x_i)=1
\right\}.
\]
$t_i^\star$ is a conditional per-call quantity, not a globally optimal joint
policy over the full trajectory. At deployment time a router observes only the
current prefix $x_i$: it cannot change earlier outputs and cannot choose future
calls before their prefixes exist.

Because per-step correctness cannot always be judged in isolation, we
approximate $V_i$ through end-to-end execution: a downgraded step is
accepted when the full trajectory still passes with the same number of
steps, which lets us infer that the replaced step was handled correctly.
The released label $\hat{t}_i$ is therefore an execution-verified
estimate of $t_i^\star$ under our fixed downgrade-and-cascade protocol
(\S\ref{sec:dataset-construction}), supplemented by a manual audit
that cross-checks this inference on a random sample of steps
(\S\ref{sec:dataset-construction}, Manual audit). A tier $t$ is accepted
when at least one model in $\mathcal M_t$ resolves the mixed-model
trajectory. We use a fixed 11-model pool spanning four tiers; the full
list and grouping are in the appendix
(Table~\ref{tab:appendix-model-pool}). Changing the pool, tier
membership, harness, or verification protocol changes $\hat{t}_i$ and
constitutes a new benchmark version.

\subsection{Corpus}

The release contains 970 step-level rows from 520 trajectory instances across
five benchmarks; per-benchmark instance and step counts (and prefix-token
medians) are summarized in Appendix Table~\ref{tab:corpus-overview}. A trajectory is one
successful end-to-end run of a strong model on the original task; each step
corresponds to one LLM call inside that trajectory, and its prefix is exactly
the messages the router would see before that call
(\S\ref{sec:dataset-construction} describes the full construction). Every
row has the shape \texttt{id}, \texttt{benchmark}, \texttt{instance\_id},
\texttt{step\_index}, \texttt{total\_steps}, \texttt{messages},
\texttt{target\_tier}, and \texttt{target\_tier\_id}. No vendor model IDs
appear in the public JSONL, only the tier label. Full row examples are in
the appendix.

Appendix Table~\ref{tab:tier-distribution} summarizes the distribution of cost-saving
opportunities across tiers. The \texttt{low} tier tests whether a router can
avoid unnecessary expensive calls while preserving task success. The two middle
tiers are transitional capability bands; routers with three output classes can
merge them into a single middle band or map through the released tier ids.
SWE-bench contributes most \texttt{high}-tier rows, signaling when a router
must stay conservative in difficult code-repair steps.

mtRAG and QMSum contribute single-call rows whose prefixes may contain
multi-turn context. BFCL contains both single-turn and multi-turn function
calling: 130 instances yield 248 rows, with 29 instances having more than one
routed step. SWE-bench and PinchBench contribute multi-step agent traces. We
include all cases because production routers encounter both single-call
long-context states and sequential agent states.

\subsection{Static evaluation without an online judge}

Static scoring is deterministic arithmetic over tier labels, trajectory
membership, and token costs; the four metrics \textsc{RowPass},
\textsc{RowExact}, \textsc{TrajPass}, and \textsc{CostSave} are defined in
\S\ref{sec:static-scoring}.

\subsection{Dynamic SWE-bench track}

The dynamic track provides a harness that supports the full 500-case
SWE-bench Verified suite; this paper reports a 100-case held-out evaluation
disjoint from any static-track SWE-bench case. At each LLM call the router
chooses a concrete model from the locked pool given the current prefix,
available models, cache state, and budget. The track reports three results:
resolve rate (the official SWE-bench predicate), realized API spend, and a
leaderboard bill that adds the same flat unresolved-instance penalty on top
of API cost for every policy (\S\ref{sec:dynamic-scoring}).
